% ============================================================================
% foreword2.tex — Perspectives by Stéphane Voisin
% ============================================================================
\forewordpageright
  {images/stephane.png}
  {Perspectives}
  {St\'{e}phane Voisin}
  {Director of Sustainable Finance Research,\\ILB \& Secretary General, Fondation PARC}

\noindent{\displayfont\fontsize{24}{14}\selectfont\color{SEGreen}When climate enters the valuation equation\par}
\vspace{6mm}

\begin{multicols}{2}\raggedcolumns
\color{SEDarkGreen}
\noindent Climate-related risk is rapidly integrating into the financial landscape, shifting from an abstract concern to a tangible factor embedded in the design of risk models across institutions. The tools used to capture this very specific nature of environmental risk must evolve accordingly. Early initiatives---such as the Dutch Central Bank's \textit{Underwater} study---marked an important step by introducing Discounted Cash Flow (DCF) analysis into climate stress testing. Yet these approaches largely treated climate shocks as external inputs applied onto otherwise unchanged financial frameworks. The Climate Cost of Computing work by SE contributes to advancing a more structural approach: a rethinking of valuation itself, where climate is no longer a scenario layered on top, but a variable within the system.

\vspace{4mm}

\noindent Infrastructure is certainly the asset class where this shift is most needed. Data centres, the silent engines of AI, are scaling at unprecedented speed, while remaining deeply exposed to physical climate constraints: heat, water stress, grid instability. Their long lifespans collide with uncertain climate trajectories, and yet their valuation still assumes a form of environmental neutrality. This paper brings that contradiction into focus.

\columnbreak

\noindent By embedding climate dynamics directly into valuation through a unified framework, it moves beyond stress testing to measure how risk propagates through costs, operations and asset value itself. It offers a bridge between climate science and financial economics, operationalising concepts like Climate Value at Risk into a tractable, decision-oriented methodology.

\vspace{4mm}

\noindent For academic research, it opens a path toward integrating physical and transition risks within standard valuation theory---transforming climate risk from a qualitative overlay into a quantifiable and actionable dimension of economic analysis.

\end{multicols}
